\documentclass[UTF8,11pt]{ctexart}
\usepackage[a4paper,margin=22mm]{geometry}
\usepackage{graphicx,enumitem,amssymb,tabularx}
\usepackage[normalem]{ulem}
\usepackage{xeCJKfntef}
\setlength{\parindent}{0pt}
\setlength{\parskip}{.6em}
\sloppy
\begin{document}
\section*{2016年6月英语六级真题(第2套)}
\subsection*{\textbf{2016} 年\textbf{6} 月大学英语六级考试真题（第二套）}


\subsection*{\textbf{Part I Writing (30 minutes)}}


For this part, you are allowed 30 minutes to write a short essay on e-learning. Try to imagine what will happen when more and more people study online instead of attending school. You are required to write at least \uline{150} words but no more than\uline{ 200} words.


\subsection*{\textbf{Section A}}


\textbf{Directions:} In this section, you will hear two long conversations. At the end of each conversation, you will hear four questions. Both the conversation and the questions will be spoken only once. After you hear a question, you must choose the best answer. from the four choices marked A), B),C) and D). Then mark the corresponding letter on Answer Sheet 1 with a single line through the centre.


\textbf{Questions 1 to 4 are based on the conversation you have just heard.}


1.


\begin{itemize}[leftmargin=2em,itemsep=.2em,topsep=.4em]\item[A.] The project the man managed at CucinTech.
\item[B.] The updating of technology at CucinTech.
\item[C.] The man's switch to a new career.
\item[D.] The restructuring of her company.
\end{itemize}

\noindent\begin{tabularx}{\linewidth}{@{}lXXXX@{}}2. & \textbf{A.} Talented personnel. & \textbf{B.} Strategic innovation. & \textbf{C.} Competitive products. & \textbf{D.} Effective promotion.\\\end{tabularx}\par

3.


\begin{itemize}[leftmargin=2em,itemsep=.2em,topsep=.4em]\item[A.] Expand the market.
\item[B.] Recruit more talents.
\item[C.] Innovate constantly.
\item[D.] Watch out for his competitors.
\end{itemize}

4.


\begin{itemize}[leftmargin=2em,itemsep=.2em,topsep=.4em]\item[A.] Possible bankruptcy.
\item[B.] Unforeseen difficulties.
\item[C.] Conflicts within the company.
\item[D.] Imitation by one's competitors.
\end{itemize}

\textbf{Questions 5 to 8 are based on the conversation you have just heard.}


5.


\begin{itemize}[leftmargin=2em,itemsep=.2em,topsep=.4em]\item[A.] The job of an interpreter.
\item[B.] The stress felt by professionals.
\item[C.] The importance of language proficiency.
\item[D.] The best way to effective communication.
\end{itemize}

\noindent\begin{tabularx}{\linewidth}{@{}lXXXX@{}}6. & \textbf{A.} Promising. & \textbf{B.} Admirable. & \textbf{C.} Rewarding. & \textbf{D.} Meaningful.\\\end{tabularx}\par

7.


\begin{itemize}[leftmargin=2em,itemsep=.2em,topsep=.4em]\item[A.] They all have a strong interest in language.
\item[B.] They all have professional qualifications.
\item[C.] They have all passed language proficiency tests.
\item[D.] They have all studied cross-cultural differences.
\end{itemize}

8.


\begin{itemize}[leftmargin=2em,itemsep=.2em,topsep=.4em]\item[A.] It requires a much larger vocabulary.
\item[B.] It attaches more importance to accuracy.
\item[C.] It is more stressful than simultaneous interpreting.
\item[D.] It puts one's long-term memory under more stress.
\end{itemize}

.


\subsection*{\textbf{Section B}}


\textbf{Directions:} In this section, you will hear two passages. At the end of each passage, you willhear three or four questions. Both the passage and the questions will be spoken only once.After you hear a question, you must choose the best answer from the four choices marked A),B), C) and D). Then mark the corresponding letter on Answer Sheet 1 with a single line throughthe centre.


\subsection*{\textbf{Passage One}}


\textbf{Questions 9 to 11 are based on the passage you have just heard.}


9.


\begin{itemize}[leftmargin=2em,itemsep=.2em,topsep=.4em]\item[A.] It might affect mothers' health.
\item[B.] It might disturb infants' sleep.
\item[C.] It might increase the risk of infants, death.
\item[D.] It might increase mothers' mental distress.
\end{itemize}

10.


\begin{itemize}[leftmargin=2em,itemsep=.2em,topsep=.4em]\item[A.] Mothers who breast-feed their babies have a harder time falling asleep.
\item[B.] Mothers who sleep with their babies need a little more sleep each night.
\item[C.] Sleeping patterns of mothers greatly affect their newborn babies' health.
\item[D.] Sleeping with infants in the same room has a negative impact on mothers.
\end{itemize}

11.


\begin{itemize}[leftmargin=2em,itemsep=.2em,topsep=.4em]\item[A.] Change their sleep patterns to adapt to their newborn babies'.
\item[B.] Sleep in the same room but not in the same bed as their babies.
\item[C.] Sleep in the same house but not in the same room as their babies.
\item[D.] Take precautions to reduce the risk of sudden infant death syndrome.
\end{itemize}

\textbf{Questions 12 to 15 are based on the passage you have just heard.}


12.


\begin{itemize}[leftmargin=2em,itemsep=.2em,topsep=.4em]\item[A.] A lot of native languages have already died out in the US.
\item[B.] The US ranks first in the number of endangered languages.
\item[C.] The efforts to preserve Indian languages have proved fruitless.
\item[D.] More money is needed to record the native languages in the US.
\end{itemize}

13.


\begin{itemize}[leftmargin=2em,itemsep=.2em,topsep=.4em]\item[A.] To set up more language schools.
\item[B.] To document endangered languages.
\item[C.] To educate native American children.
\item[D.] To revitalise America's native languages.
\end{itemize}

14.


\begin{itemize}[leftmargin=2em,itemsep=.2em,topsep=.4em]\item[A.] The US govemment's policy of Americanising Indian children.
\item[B.] The failure of American Indian languages to gain an official status.
\item[C.] The US government's unwillingness to spend money educating Indians.
\item[D.] The long-time isolation of American Indians from the outside world.
\end{itemize}

15.


\begin{itemize}[leftmargin=2em,itemsep=.2em,topsep=.4em]\item[A.] It is being utilised to teach native languages.
\item[B.] It tells traditional stories during family time.
\item[C.] It speeds up the extinction of native languages.
\item[D.] It is widely used in language immersion schools.
\end{itemize}

\subsection*{\textbf{Section C}}


\textbf{Directions:} In this section, you will hear three recordings of lectures or talks followed by threeor four questions. The recordings will be played only once. After you hear a question, you mustchoose the best answer from the four choices marked A),B),C) and D). Then mark thecorresponding letter on Answer Sheet 1 with a single line through the centre. \textbf{Recording One}


\textbf{Questions 16 to 18 are based on the recording you have just heard.}


16.


\begin{itemize}[leftmargin=2em,itemsep=.2em,topsep=.4em]\item[A.] It pays them up to half of their previous wages while they look for work.
\item[B.] It covers their mortgage payments and medical expenses for 99 weeks.
\item[C.] It pays their living expenses until they find employment again.
\item[D.] It provides them with the basic necessities of everyday life.
\end{itemize}

17.


\begin{itemize}[leftmargin=2em,itemsep=.2em,topsep=.4em]\item[A.] Creating jobs for the huge army of unemployed workers.
\item[B.] Providing training and guidance for unemployed workers.
\item[C.] Convincing local lawmakers to extend unemployment benefits.
\end{itemize}

\begin{itemize}[leftmargin=2em,itemsep=.2em,topsep=.4em]\item[D.] Raising funds to help those having no unemployment insurance.
\end{itemize}

18.


\begin{itemize}[leftmargin=2em,itemsep=.2em,topsep=.4em]\item[A.] To offer them loans they need to start their own businesses.
\item[B.] To allow them to postpone their monthly mortgage payments.
\item[C.] To create more jobs by encouraging private investments in local companies.
\item[D.] To encourage big businesses to hire back workers with government subsidies.
\end{itemize}

\textbf{Questions 19 to 22 are based on the recording you have just heard.}


19.


\begin{itemize}[leftmargin=2em,itemsep=.2em,topsep=.4em]\item[A.] They measured the depths of sea water.
\item[B.] They analyzed the water content.
\item[C.] They explored the ocean floor.
\item[D.] They investigated the ice.
\end{itemize}

20.


\begin{itemize}[leftmargin=2em,itemsep=.2em,topsep=.4em]\item[A.] Eighty percent of the ice disappears in summer time.
\item[B.] Most of the ice was accumulated over the past centuries.
\item[C.] The ice ensures the survival of many endangered species.
\item[D.] The ice decrease is more evident than previously thought.
\end{itemize}

21.


\begin{itemize}[leftmargin=2em,itemsep=.2em,topsep=.4em]\item[A.] Arctic ice is a major source of the world's fresh water.
\item[B.] The melting Arctic ice has drowned many coastal cities.
\item[C.] The decline of Arctic ice is irreversible.
\item[D.] Arctic ice is essential to human survival.
\end{itemize}

22.


\begin{itemize}[leftmargin=2em,itemsep=.2em,topsep=.4em]\item[A.] It will do a lot of harm to mankind.
\item[B.] There is no easy way to understand it.
\item[C.] It will advance nuclear technology.
\item[D.] There is no easy technological solution to it.
\end{itemize}

\textbf{Questions 23 to 25 are based on the recording you have just heard.}


23.


\begin{itemize}[leftmargin=2em,itemsep=.2em,topsep=.4em]\item[A.] The reason why New Zealand children seem to have better self-control.
\item[B.] The relation between children's self-control and their future success.
\item[C.] The health problems of children raised by a single parent.
\item[D.] The deciding factor in children's academic performance.
\end{itemize}

24.


\begin{itemize}[leftmargin=2em,itemsep=.2em,topsep=.4em]\item[A.] Children raised by single parents will have a hard time in their thirties.
\item[B.] Those with a criminal record mostly come from single parent families.
\item[C.] Parents must learn to exercise self-control in front of their children.
\item[D.] Lack of self-control in parents is a disadvantage for their children.
\end{itemize}

25.


\begin{itemize}[leftmargin=2em,itemsep=.2em,topsep=.4em]\item[A.] Self-control can be improved through education.
\item[B.] Self-control can improve one's financial situation.
\item[C.] Self-control problems may be detected early in children.
\item[D.] Self-control problems will diminish as one grows up.
\end{itemize}

\subsection*{\textbf{Directions:}}


\subsection*{\textbf{Part III Reading comprehension (40 minutes)}}


\subsection*{\textbf{Section A}}


\textbf{Directions:} In this section, there is a passage with ten blanks. You are required to select one word for each blank from a list of choices given in a word bank following the passage. Read the passage through carefully before making your choices. Each choice in the bank is identified by a letter. Please mark the corresponding letter for each item on Answer Sheet 2 with a single line through the center. You may not use any of the words in the bank more than once.


The robotics revolution is set to bring humans face to face with an old fear—man-made creations as smart and capable as we are but without a moral compass. As robots take on ever more complex roles, the question naturally\uline{ 26 }: Who will be responsible when they do something wrong? Manufacturers? Users? Software writers? The answer depends on the robot.


Robots already save us time, money and energy. In the future, they will improve our health care, social welfare and standard of living. The\uline{ 27 }of computational power and engineering advances will\uline{ 28 }enable lower-cost in- home care for the disabled,\uline{ 29} use of driverless cars that may reduce drunk- and distracted-driving accidents and countless home and service-industry uses for robots, from street cleaning to food preparation.


But there are\uline{ 30} to be problems. Robot cars will crash. A drone ( 遥控飞行器 ) operator will\uline{ 31 }someone’s privacy. A robotic lawn mower will run over a neighbor’s cat. Juries sympathetic to the\uline{ 32 }of machines will punish entrepreneurs with company-crushing \uline{33 }and damages. What should governments do to protect people while\uline{ 34} \uline{s}pace for innovation?


Big, complicated systems on which much public safety depends, like driverless cars, should be built, \uline{35 }and sold by manufacturers who take responsibility for ensuring safety and are liable for accidents. Governments should set safety requirements and then let insurers price the risk of the robots based on the manufacturer’s driving record, not the passenger’s.


\begin{itemize}[leftmargin=2em,itemsep=.2em,topsep=.4em]\item[A.] arises
\item[B.] ascends
\item[C.] bound
\item[D.] combination
\item[E.] definite
\item[F.] eventually
\item[G.] interfere
\item[H.] invade
\item[I.] manifesting
\item[J.] penalties
\item[K.] preserving
\item[L.] programmed
\item[M.] proximately
\item[N.] victims
\item[O.] widespread
\end{itemize}

\subsection*{\textbf{Section B}}


\textbf{Directions:} In this section, you are going to read a passage with ten statements attached to it. Each statement contains information given in one of the paragraphs. Identify the paragraph from which the information is derived. You may choose a paragraph more than once. Each paragraph is marked with a letter. Answer the questions by marking the corresponding letter on Answer Sheet 2. Reform and Medical Costs


\begin{itemize}[leftmargin=2em,itemsep=.2em,topsep=.4em]\item[A.] Americans are deeply concerned about the relentless rise in health care costs and health insurance premiums.
\end{itemize}

They need to know if reform will help solve the problem. The answer is that no one has an easy fix for rising medical costs. The fundamental fix—reshaping how care is delivered and how doctors are paid in a wasteful, abnormal system—is likely to be achieved only through trial and error and incremental ( 渐进的 ) gains.


\begin{itemize}[leftmargin=2em,itemsep=.2em,topsep=.4em]\item[B.] The good news is that a bill just approved by the House and a bill approved by the Senate Finance Committee
\end{itemize}

would implement or test many reforms that should help slow the rise in medical costs over the long term. As a report in The New England Journal of Medicine concluded, “Pretty much every proposed innovation found in the health policy Iiterature these days is contained in these measures.”


\begin{itemize}[leftmargin=2em,itemsep=.2em,topsep=.4em]\item[C.] Medical spending, which typically rises faster than wages and the overall economy, is propelled by two things:
\end{itemize}

the high prices charged for medical services in this country and the volume of unnecessary care delivered by doctors and hospitals, which often perform a lot more tests and treatments than a patient really needs.


\begin{itemize}[leftmargin=2em,itemsep=.2em,topsep=.4em]\item[D.] Here are some of the important proposals in the House and Senate bills to try to address those problems, and
\end{itemize}

why it is hard to know how well they will work.


\begin{itemize}[leftmargin=2em,itemsep=.2em,topsep=.4em]\item[E.] Both bills would reduce the rate of growth in annual Medicare payments to hospitals, nursing homes and other
\end{itemize}

providers by amounts comparable to the productivity savings routinely made in other industries with the help of new technologies and new ways to organize work. This proposal could save Medicare more than \$100 billion over the next decade. If private plans demanded similar productivity savings from providers, and refused to let providers shift additional costs to them, the savings could be much larger. Critics say Congress will give in to lobbyists and let inefficient providers off the hook (放过). That is far less likely to happen if Congress also adopts strong “pay-go” rules requiring that any increase in payments to providers be offset by new taxes or budget cuts.


\begin{itemize}[leftmargin=2em,itemsep=.2em,topsep=.4em]\item[F.] The Senate Finance bill would impose an excise tax ( 消费税 ) on health insurance plans that cost more than
\end{itemize}

\$8,000 for an individual or \$21,000 for a family. It would most likely cause insurers to redesign plans to fall beneath the threshold. Enrollees would have to pay more money for many services out of their own pockets, and that would encourage them to think twice about whether an expensive or redundant test was worth it. Economists project that most employers would shift money from expensive health benefits into wages, The House bill has no similar tax. The final legislation should.


\begin{itemize}[leftmargin=2em,itemsep=.2em,topsep=.4em]\item[G.] Any doctor who has wrestled with multiple forms from different insurers, or patients who have tried to
\end{itemize}

understand their own parade of statements, know that simplification ought to save money. When the health insurance industry was still cooperating in reform efforts, its trade group offered to provide standardized forms for automated processing. It estimated that step would save hundreds of billions of dollars over the next decade. The bills would lock that pledge into law.


\begin{itemize}[leftmargin=2em,itemsep=.2em,topsep=.4em]\item[H.] The stimulus package provided money to convert the inefficient, paper-driven medical system to electronic
\end{itemize}

records that can be easily viewed and transmitted. This requires open investments to help doctors convert. In time it should help restrain costs by eliminating redundant tests, preventing drug interactions, and helping doctors find the best treatments.


\begin{itemize}[leftmargin=2em,itemsep=.2em,topsep=.4em]\item[I.] Virtually all experts agree that the fee-for-service system—doctors are rewarded for the quantity of care rather
\end{itemize}

than its quality or effectiveness—is a primary reason that the cost of care is so high. Most agree that the solution is to push doctors to accept fixed payments to care for a particular illness or for a patient’s needs over a year. No one knows how to make that happen quickly. The bills in both houses would start pilot projects within Medicare. They include such measures as accountable care organizations to take charge of a patient’s needs with an eye on both cost and quality, and chronic disease management to make sure the seriously ill, who are responsible for the bulk of all health care costs, are treated properly. For the most part, these experiments rely on incentive payments to get doctors to try them.


\begin{itemize}[leftmargin=2em,itemsep=.2em,topsep=.4em]\item[J.] Testing innovations do no good unless the good experiments are identified and expanded and the bad ones are
\end{itemize}

dropped. The Senate bill would create an independent commission to monitor the pilot programs and recommend changes in Medicare’s payment policies to urge providers to adopt reforms that work. The changes would have to be approved or rejected as a whole by Congress, making it hard for narrow-interest lobbies to bend lawmakers to their will.


\begin{itemize}[leftmargin=2em,itemsep=.2em,topsep=.4em]\item[K.] The bills in both chambers would create health insurance exchanges on which small businesses and individuals
\end{itemize}

could choose from an array of private plans and possibly a public option. All the plans would have to provide standard benefit packages that would be easy to compare. To get access to millions of new customers, insurers would have a strong incentive to sell on the exchange. And the head-to-head competition might give them a strong incentive to lower their prices, perhaps by accepting slimmer profit margins or demanding better deals from providers.


\begin{itemize}[leftmargin=2em,itemsep=.2em,topsep=.4em]\item[L.] The final legislation might throw a public plan into the competition, but thanks to the fierce opposition of the
\end{itemize}

insurance industry and Republican critics, it might not save much money. The one in the House bill would have to negotiate rates with providers, rather than using Medicare rates, as many reformers wanted.


\begin{itemize}[leftmargin=2em,itemsep=.2em,topsep=.4em]\item[M.] The president’s stimulus package is pumping money into research to compare how well various treatments
\end{itemize}

work. Is surgery, radiation or careful monitoring best for prostate ( 前列腺) cancer? Is the latest and most expensive cholesterol-lowering drug any better than its common competitors? The pending bills would spend additional money to accelerate this effort.


\begin{itemize}[leftmargin=2em,itemsep=.2em,topsep=.4em]\item[N.] Critics have charged that this sensible idea would lead to rationing of care. (That would be true only if you
\end{itemize}

believed that patients should have an unrestrained right to treatments proven to be inferior.) As a result, the bills do not require, as they should, that the results of these studies be used to set payment rates in Medicare.


\begin{itemize}[leftmargin=2em,itemsep=.2em,topsep=.4em]\item[O.] Congress needs to find the courage to allow Medicare to pay preferentially for treatments proven to be superior.
\end{itemize}

Sometimes the best treatment might be the most expensive. But overall, we suspect that spending would come down through elimination of a lot of unnecessary or even dangerous tests and treatments.


[P] The House bill would authorize the secretary of health and human services to negotiate drug prices in Medicare and Medicaid. Some authoritative analysts doubt that the secretary would get better deals than private insurers already get. We believe negotiation could work. It does in other countries.


[Q] Missing from these bills is any serious attempt to rein in malpractice costs. Malpractice awards do drive up insurance premiums for doctors in high-risk specialties, and there is some evidence that doctors engage in “defensive medicine” by performing tests and treatments primarily to prove they are not negligent should they get sued.


36. With a tax imposed on expensive health insurance plans, most employers will likely transfer money from health


expenses into wages.


37. Changes in policy would be approved or rejected as a whole so that lobbyists would find it hard to influence


lawmakers.


38. It is not easy to curb the rising medical costs in America.


39. Standardization of forms for automatic processing will save a lot of medical expenses.


40. Republicans and the insurance industry are strongly opposed to the creation of a public insurance plan.


41. Conversion of paper to electronic medical records will help eliminate redundant tests and prevent drug


interactions.


42. The high cost of medical services and unnecessary tests and treatments have driven up medical expenses.


43. One main factor that has driven up medical expenses is that doctors are compensated for the amount of care rather


than its effect.


44. Contrary to analysts’ doubts, the author believes drug prices may be lowered through negotiation.


45. Fair competition might create a strong incentive for insurers to charge less.


\subsection*{\textbf{Section C}}


\textbf{Directions:} There are two passages in this section. Each passage is followed by some questions or unfinished statements. For each of them there are four choices marked A), B), C) and D).You should decide on the best choice and mark the corresponding letter on Answer Sheet 2 with a single line through the center.


\subsection*{\textbf{Passage One}}


\textbf{Questions 46 to 50 are based on the following passage.}


Facing water shortages and escalating fertilizer costs, farmers in developing countries are using raw sewage ( 下水道污水 ) to irrigate and fertilize nearly 49 million acres of cropland, according to a new report—and it may not be a bad thing.


While the practice carries serious health risks for many, those dangers are outweighed by the social and economic gains for poor urban farmers and consumers who need affordable food. “There is a large potential for wastewater agriculture to both help and hurt great numbers of urban consumers,” said Liqa Raschid-Sally, who led the study.


The report focused on poor urban areas, where farms in or near cities supply relatively inexpensive food. Most of these operations draw irrigation water from local rivers or lakes. Unlike developed cities, however, these areas lack advanced water-treatment facilities, and rivers effectively become sewers ( 下水道 ).


When this water is used for agricultural irrigation, farmers risk absorbing disease-causing becteria, as do consumers who eat the produce raw and unwashed. Nearly 2.2 million people die a year because of diarrhea-related ( 与腹泻相关的) diseases, according to WHO statistics. More than 80\% of those cases can be attributed to contact with contaminated water and a lack of proper sanitation. But Pay Drechsel, an environmental scientist, argues that the social and economic benefits of using untreated human waste to grow food outweigh the health risks.


Those dangers can be addressed with farmer and consumer education, he said, while the free water and nutrients from human waste can help urban farmers in developing countries to escape poverty.


Agriculture is a water-intensive business, accounting for nearly 70\% of global fresh water consumption.


In poor, dry regions, untreated wastewater is the only viable irrigation source to keep farmers in business. In some cases, water is so scarce that farmers break open sewage pipes transporting waste to local rivers. Irrigation is the primary agricultural use of human waste in the developing world. But frequently untreated human waste harvested from lavatories is delivered to farms and spread as fertilizer.


In most cases, the human waste is used on grain crops, which are eventually cooked, minimizing the risk of transmitting water-borne diseases. With fertilizer prices jumping nearly 50\% per metric ton over the last year in some places, human waste is an attractive, and often necessary, alternative.


In cases where sewage mud is used, expensive chemical fertilizer use can be avoided. The mud contains the same critical nutrients.


“Overly strict standards often fail,” James Bartram, a WHO water-health expert, said. “We need to accept that fact across much of the planet, so waste with little or no treatment will be used in agriculture for good reason.”


46. What does the author say about the use of raw sewage for farming?


\begin{itemize}[leftmargin=2em,itemsep=.2em,topsep=.4em]\item[A.] Its risks cannot be overestimated.
\item[B.] It should be forbidden altogether.
\item[C.] Its benefits outweigh the hazards involved.
\item[D.] It is polluting millions of acres of cropland.
\end{itemize}

47. What is the main problem caused by the use of wastewater for irrigation?


\begin{itemize}[leftmargin=2em,itemsep=.2em,topsep=.4em]\item[A.] Rivers and lakes nearby will gradually become contaminated.
\item[B.] It will drive producers of chemical fertilizers out of business.
\item[C.] Farmers and consumers may be affected by harmful bacteria.
\item[D.] It will make the farm produce less competitive on the market.
\end{itemize}

48. What is environmental scientist Pay Drechsel’s attitude towards the use of untreated human waste in agriculture?


\begin{itemize}[leftmargin=2em,itemsep=.2em,topsep=.4em]\item[A.] Favorable.
\item[B.] Skeptical.
\item[C.] Indifferent.
\item[D.] Responsible.
\end{itemize}

49. What does Pay Dreschsel think of the risks involved in using untreated human waste for farming?


\begin{itemize}[leftmargin=2em,itemsep=.2em,topsep=.4em]\item[A.] They have been somewhat exaggerated.
\item[B.] They can be dealt with through education.
\item[C.] They will be minimized with new technology.
\item[D.] They can be addressed by improved sanitation.
\end{itemize}

50. What do we learn about James Bartram’s position on the use of human waste for farming?


\begin{itemize}[leftmargin=2em,itemsep=.2em,topsep=.4em]\item[A.] He echoes Pay Drechsel’s opinion on the issue.
\item[B.] He challenges Liqa Raschid-Sally’s conclusion
\item[C.] He thinks it the only way out of the current food crisis.
\item[D.] He deems it indispensable for combating global poverty.
\end{itemize}

\subsection*{\textbf{Passage Two}}


\textbf{Questions 51 to 55 are based on the following passage.}


These days, nobody needs to cook. Families graze on high-cholesterol take-aways and microwaved ready-meals. Cooking is an occasional hobby and a vehicle for celebrity chefs. Which makes it odd that the kitchen has become the heart of the modern house: what the great hall was to the medieval castle, the kitchen is to the 21st-century home. The money spent on kitchens has risen with their status. In America the kitchen market is now worth \$170 billion, 5 times the country’s film industry. In the year to August 2007, IKEA, a Swedish furniture chain, sold over one million kitchens worldwide. The average budget for a “major” kitchen overhaul in 2006, calculates Remodeling magazine, was a staggering \$54,000; even a “minor” improvement cost on average \$18,000.


Exclusivity, more familiar in the world of high fashion, has reached the kitchen: Robinson \& Cornish, a British manufacturer of custom-made kitchens, offers a Georgian-style one which would cost £145,000–155,000— excluding building, plumbing and electrical work. Its big selling point is that nobody else will have it: “You won’t see this kitchen anywhere else in the world.”


The elevation of the room that once belonged only to the servants to that of design showcase for the modern family tells the story of a century of social change. Right into the early 20th century, kitchens were smoky, noisy places, generally located underground, or to the back of the house, and as far from living space as possible. That was as it should be: kitchens were for servants, and the aspiring middle classes wanted nothing to do with them.


But as the working classes prospered and the servant shortage set in, housekeeping became a matter of interest to the educated classes. One of the pioneers of a radical new way of thinking about the kitchen was Catharine Esther Beecher, sister of Harriet Beecher Stowe. In American Woman’s Home, published in 1869, the Beecher sisters recommended a scientific approach to household management, designed to enhance the efficiency of a woman’s work and promote order.


Many contemporary ideas about kitchen design can be traced back to another American, Christine Frederick, who set about enhancing the efficiency of the housewife. Her 1919 work, Household Engineering: Scientific Management in the Home, was based on detailed observation of a housewife’s daily routine. She borrowed the principle of efficiency on the factory floor and applied it to domestic tasks on the kitchen floor. Frederick’s central idea, that “stove, sink and kitchen table must be placed in such a relation that useless steps are avoided entirely,” inspired the first fully fitted kitchen, designed in the 1920s by Margarete Schütter-Lihotsky. It was a modernist triumph, and many elements remain central features of today’s kitchen.


51. What does the author say about the kitchen of today?


\begin{itemize}[leftmargin=2em,itemsep=.2em,topsep=.4em]\item[A.] It is where housewives display their cooking skills.
\item[B.] It is where the family entertains important guests.
\item[C.] It has become something odd in a modern house.
\item[D.] It is regarded as the center of a modern home.
\end{itemize}

52. Why does the Georgian-style kitchen sell at a very high price?


\begin{itemize}[leftmargin=2em,itemsep=.2em,topsep=.4em]\item[A.] It is believed to have tremendous artistic value.
\item[B.] No duplicate is to be found in any other place.
\item[C.] It is manufactured by a famous British company.
\item[D.] No other manufacturer can produce anything like it.
\end{itemize}

53. What does the change in the status of the kitchen reflect?


\begin{itemize}[leftmargin=2em,itemsep=.2em,topsep=.4em]\item[A.] Improved living conditions.
\item[B.] Women’s elevated status.
\item[C.] Technological progress.
\item[D.] Social change.
\end{itemize}

54. What was the Beecher sisters’ idea of a kitchen?


\begin{itemize}[leftmargin=2em,itemsep=.2em,topsep=.4em]\item[A.] A place where women could work more efficiently.
\item[B.] A place where high technology could be applied.
\item[C.] A place of interest to the educated people.
\item[D.] A place to experiment with new ideas.
\end{itemize}

55. What do we learn about today’s kitchen?


\begin{itemize}[leftmargin=2em,itemsep=.2em,topsep=.4em]\item[A.] It represents the rapid technological advance in people’s daily life.
\item[B.] Many of its central features are no different from those of the 1920s.
\item[C.] It has been transformed beyond recognition.
\item[D.] Many of its functions have changed greatly.
\end{itemize}

Part IVTranslation (30 minutes) Directions: For this part, you are allowed 30 minutes to translate a passage from Chinese into English. You should write your answer on Answer Sheet 2. 旗袍（qipao）是一种雅致的中国服装，源于中国的满族（Manchu Nationality）。在清代，旗袍是王室女性穿着的宽松长袍。上世纪20 年代，受西方服饰影响，旗袍发生了一些变化。袖口（cuffs）变窄，袍身变短。这些变化使女性美得以充分展现。如今，旗袍经常出现在世界级的时装秀上。中国女性出席重要社交聚会时，旗袍往往是她们的首选。很多中国新娘也会选择旗袍作为结婚礼服。一些有影响的人士甚至建议将旗袍作为中国女性的民族服饰。

\end{document}